%==============================================================
%  lecons/analyse/tfa.tex — Théorème fondamental de l'analyse
%==============================================================

\lecontitle{Théorème fondamental de l'analyse}{Analyse réelle — Intégration et primitives}

%=============================================================
%  § 1 — ÉNONCÉ
%=============================================================
\secnum{navyblue}{1}{Énoncé}

\begin{tcolorbox}[thmbox]
  \textbf{Théorème fondamental de l'analyse.}\enspace%
  \index{théorème fondamental de l'analyse}
  Soit $f : [a,b] \to \mathbb{R}$ \textbf{continue}.

  \medskip
  \textbf{Partie I — L'intégrale définit une primitive.}\enspace%
  \textit{La fonction $F$ définie par}
  \[
    F(x) = \int_a^x f(t)\,dt
  \]
  \textit{est l'unique primitive de $f$ sur $[a,b]$ qui s'annule en $a$. En particulier,
  $F$ est dérivable sur $[a,b]$ et $F'(x) = f(x)$ pour tout $x \in [a,b]$.}

  \medskip
  \textbf{Partie II — Formule de Newton–Leibniz.}\enspace%
  \textit{Si $G$ est une primitive quelconque de $f$ sur $[a,b]$, alors}
  \[
    \int_a^b f(t)\,dt = G(b) - G(a) =: \bigl[G(t)\bigr]_a^b.
  \]
\end{tcolorbox}

\rubriq{Signification}
La Partie~I dit que l'intégration et la dérivation sont des opérations inverses :
dériver l'intégrale de $f$ redonne $f$. La Partie~II donne le calcul effectif :
pour évaluer une intégrale, il suffit de trouver une primitive.
Ensemble, elles établissent l'équivalence entre le calcul différentiel et le calcul
intégral, deux théories développées indépendamment (Newton, Leibniz).

\decsep{}

%=============================================================
%  § 2 — DÉMONSTRATION
%=============================================================
\secnum{navyblue}{2}{Démonstration}

\rubriq{Preuve de la Partie I}

\begin{mdframed}[style=proofbar]
  Soit $x \in [a,b]$ et $h \neq 0$ assez petit. On a
  \[
    F(x+h) - F(x) = \int_x^{x+h} f(t)\,dt.
  \]
  On veut montrer que $\dfrac{F(x+h)-F(x)}{h} \to f(x)$ quand $h\to 0$.
  On écrit
  \[
    \frac{F(x+h)-F(x)}{h} - f(x)
    = \frac{1}{h}\int_x^{x+h}\bigl(f(t)-f(x)\bigr)\,dt.
  \]
  Soit $\varepsilon > 0$. Par continuité de $f$ en $x$, il existe $\delta > 0$ tel que
  $|t-x| \leq |h| < \delta \Rightarrow |f(t)-f(x)| < \varepsilon$. Alors
  \[
    \left|\frac{F(x+h)-F(x)}{h} - f(x)\right|
    \leq \frac{1}{|h|}\int_x^{x+h}|f(t)-f(x)|\,dt \leq \varepsilon.
  \]
  Donc $F'(x) = f(x)$. L'unicité vient du fait que deux primitives diffèrent
  d'une constante, et $F(a) = 0$ fixe celle-ci. \hfill$\blacksquare$
\end{mdframed}

\rubriq{Preuve de la Partie II}

\begin{mdframed}[style=proofbar]
  Soit $G$ une primitive de $f$. Alors $(G-F)' = f - f = 0$ sur $[a,b]$,
  donc $G - F$ est constante (corollaire du TAF\index{accroissements finis (théorème des)}).
  En $a$ : $G(a) - F(a) = G(a) - 0 = G(a)$. Donc pour tout $x$ :
  \[
    G(x) - F(x) = G(a),
    \qquad\text{soit}\quad
    F(x) = G(x) - G(a).
  \]
  En $x = b$ : $\int_a^b f(t)\,dt = F(b) = G(b) - G(a)$. \hfill$\blacksquare$
\end{mdframed}

\decsep{}

%=============================================================
%  § 3 — COROLLAIRES
%=============================================================
\secnum{navyblue}{3}{Corollaires importants}

\begin{tcolorbox}[thmbox]
  Soient $f, g$ continues et $G, H$ des primitives respectives.
  \begin{enumerate}
    \item \textbf{Intégration par parties :}\index{intégration par parties}
    \[
      \int_a^b f(t)\,g'(t)\,dt = \bigl[f(t)g(t)\bigr]_a^b - \int_a^b f'(t)\,g(t)\,dt.
    \]
    \item \textbf{Changement de variable :}\index{changement de variable}
    Si $\varphi : [\alpha,\beta]\to[a,b]$ est $\mathcal{C}^1$, $\varphi(\alpha)=a$, $\varphi(\beta)=b$ :
    \[
      \int_a^b f(t)\,dt = \int_\alpha^\beta f\bigl(\varphi(u)\bigr)\,\varphi'(u)\,du.
    \]
    \item \textbf{Dérivation sous le signe intégral :}
    Si $f(x,t)$ et $\partial_x f(x,t)$ sont continues, alors
    \[
      \frac{d}{dx}\int_a^b f(x,t)\,dt = \int_a^b \partial_x f(x,t)\,dt.
    \]
    \item \textbf{Paramètre aux bornes :}
    $\displaystyle\frac{d}{dx}\int_a^{u(x)} f(t)\,dt = f\bigl(u(x)\bigr)\cdot u'(x)$.
  \end{enumerate}
\end{tcolorbox}

\rubriq{Preuve de l'intégration par parties}

\begin{mdframed}[style=proofbar]
  $(fg)' = f'g + fg'$ s'intègre sur $[a,b]$ via Newton-Leibniz :
  $[fg]_a^b = \int_a^b f'g + \int_a^b fg'$. \hfill$\blacksquare$
\end{mdframed}

\decsep{}

%=============================================================
%  § 4 — APPLICATIONS, CONTRE-EXEMPLES, EXERCICES
%=============================================================
\secnum{exgreen}{4}{Applications, contre-exemples et exercices}

\noindent{\bfseries\color{navyblue} Applications.}

\begin{mdframed}[style=exbar]
  \textbf{1. Calcul direct.}\enspace%
  $\displaystyle\int_0^1 t^n\,dt = \Bigl[\frac{t^{n+1}}{n+1}\Bigr]_0^1 = \frac{1}{n+1}$.

  \smallskip\noindent
  \textbf{2. Formule de Taylor avec reste intégral.}\enspace%
  \index{Taylor (reste intégral)}
  Par applications successives de la Partie~I :
  \[
    f(b) = f(a) + f'(a)(b-a) + \cdots + \frac{f^{(n)}(a)}{n!}(b-a)^n
    + \int_a^b \frac{(b-t)^n}{n!}f^{(n+1)}(t)\,dt.
  \]

  \smallskip\noindent
  \textbf{3. Sommes de Riemann.}\enspace%
  Pour $f$ continue, $\dfrac{1}{n}\displaystyle\sum_{k=0}^{n-1}f\!\left(\frac{k}{n}\right) \to \int_0^1 f(t)\,dt$.
\end{mdframed}

\vspace{6pt}
\noindent{\bfseries\color{counterred} Pièges et contre-exemples.}

\begin{mdframed}[style=cexbar]
  \textbf{La continuité est indispensable pour la Partie~I.}\enspace%
  Soit $f = \mathbf{1}_{[1/2,1]}$ sur $[0,1]$ (indicatrice, discontinue en $1/2$).
  $F(x) = \int_0^x f$ est continue et affine par morceaux, mais $F'(1/2)$ n'existe pas.

  \smallskip\noindent
  \textbf{Toute fonction dérivable n'est pas une primitive de sa dérivée
  au sens de Riemann.}\enspace%
  Si $f'$ n'est pas intégrable au sens de Riemann, la formule de Newton–Leibniz
  ne s'applique pas directement (elle est valide pour l'intégrale de Lebesgue).

  \smallskip\noindent
  \textbf{$\int_a^b f = 0$ n'implique pas $f = 0$.}\enspace%
  $\int_0^{2\pi}\sin t\,dt = 0$, mais $\sin \not\equiv 0$.
  En revanche, si $f \geq 0$ et $\int f = 0$, alors $f = 0$ (par continuité).
\end{mdframed}

\vspace{6pt}
\exolabel

\begin{mdframed}[style=exobar]
  \begin{enumerate}
    \item Calculer $\displaystyle\frac{d}{dx}\int_x^{x^2} e^{-t^2}\,dt$.

    \item Montrer que $\displaystyle f(x) = \int_0^x (x-t)\sin(t^2)\,dt$
          est de classe $\mathcal{C}^2$ et calculer $f''$.

    \item Soit $f$ continue sur $[0,1]$. Montrer que
          $\displaystyle\lim_{n\to+\infty} n\int_0^1 t^n f(t)\,dt = f(1)$.

    \item Déduire la formule de Taylor-Lagrange du reste intégral en majorant
          \[\Bigl|\int_a^b \tfrac{(b-t)^n}{n!}f^{(n+1)}(t)\,dt\Bigr| \leq \tfrac{M}{(n+1)!}|b-a|^{n+1}.\]

    \item Calculer $\displaystyle\int_0^{\pi/2}\frac{dt}{1+\tan^n(t)}$ pour $n > 0$.
          (\textit{Hint :} changement $t \mapsto \pi/2 - t$.)
  \end{enumerate}
\end{mdframed}

\leconfooter{I.\ Newton (1666) — G.\ W.\ Leibniz (1675) — A.-L.\ Cauchy (1823)}
